\section{The parity contradiction}\label{sec:parity}

We now prove Theorem~\ref{thm:parity}.  Reference qubits retain the
classical input labels while the computation may freely change all
its own registers.  Their signed acceptance expectation will be
bounded away from zero by correctness and forced to zero by
Corollary~\ref{cor:product-tail}.

\begin{proof}[Proof of Theorem~\ref{thm:parity}]
Fix positive integers $d,k,C$ and $0<\varepsilon\le1/2$, where $k$ is
now the resource exponent in the theorem. Suppose, toward a contradiction, that an unbounded
set of input lengths $n$ admits a depth-at-most-$d$ circuit $U$ with
at most $C(n+1)^k$ total qubits,
whose success probability is at least $1/2+\varepsilon$ on every input. Choose
one such circuit at each member of an increasing sequence of lengths.
Write $a$ for its number of ancillary qubits and $I$ for its $n$
input qubits.

Attach a register $R$ of $n$ reference qubits, initially zero, and
let $\operatorname{COPY}$ be the layer of disjoint CNOTs from input
qubit $I_i$ to reference qubit $R_i$.  Our normalized product starting
state is
\[
 \psi_0=|+\rangle_I^{\otimes n}
          \otimes|0\rangle^{\otimes a}
          \otimes|0\rangle_R^{\otimes n},
 \qquad |+\rangle=\frac{|0\rangle+|1\rangle}{\sqrt2}.
\]
The enlarged circuit $(U\otimes I_R)\operatorname{COPY}$ has
$m=2n+a$ qubits.  Splitting the original layers into the two analytic
types of Section~\ref{sec:framework} gives at most $t=2d+1$ layers,
including the copying layer.  In particular the layer count is
bounded, $m\to\infty$, and, since $k\ge1$,
\begin{equation}\label{eq:parity-resources}
 m=2n+a\le(C+2)(n+1)^k.
\end{equation}

Apply Proposition~\ref{prop:pruning} to this circuit and $\psi_0$,
and call the pruned circuit $W$.  Every copying gate is retained for
all sufficiently large $m$.  Indeed its reflection projection is
\[
 P_g=|1\rangle\langle1|_{I_i}
       \otimes|-\rangle\langle-|_{R_i},
 \qquad |-\rangle=\frac{|0\rangle-|1\rangle}{\sqrt2},
\]
and its probability on $\psi_0$ is
$|\langle1|+\rangle|^2|\langle-|0\rangle|^2=1/4\ge m^{-6}$
eventually.  All later gates avoid $R$, so we can write
\[
 W=(\widetilde U\otimes I_R)\operatorname{COPY}
\]
for a circuit $\widetilde U$ on the original qubits.  Let
\[
 \Phi=(U\otimes I_R)\operatorname{COPY}\psi_0,
 \qquad \Omega=W\psi_0.
\]
Both states are normalized, the trajectory to $\Omega$ is regular,
and \eqref{eq:pruning-error} gives
$\norm{\Phi-\Omega}\le2t m^{-2}$.

Let $P_B$ project the designated answer qubit onto $|1\rangle$, and
let $Z_R=\prod_{i=1}^n Z_{R_i}$.  The answer qubit belongs to the
original circuit and may be one of its input qubits; in either case
$P_B$ acts outside $R$ and commutes with $Z_R$.  Set
$p_x=\Pr_U(1\mid x)$.  The original coherent output is
\[
 \Phi=2^{-n/2}\sum_{x\in\{0,1\}^n}
        U\bigl(|x\rangle_I\otimes|0\rangle^{\otimes a}\bigr)
        \otimes|x\rangle_R.
\]
The reference states are orthogonal, so
\begin{equation}\label{eq:parity-correlation}
 \ip{\Phi}{Z_RP_B\Phi}
 =2^{-n}\sum_x(-1)^{|x|}p_x
 \le\frac12\!\left(\left(\frac12-\varepsilon\right)
                  -\left(\frac12+\varepsilon\right)\right)
 =-\varepsilon.
\end{equation}
Here $p_x\le1/2-\varepsilon$ on even-parity inputs and
$p_x\ge1/2+\varepsilon$ on
odd-parity inputs, and there are equally many of each for $n\ge1$.

The self-adjoint observable $Z_RP_B$ has norm at most one.  For any
two unit vectors its expectation changes by at most twice their
distance, by expanding the difference in its two arguments.  Hence
\begin{equation}\label{eq:pruned-correlation}
 \ip{\Omega}{Z_RP_B\Omega}
 \le-\varepsilon+4t m^{-2}.
\end{equation}
Thus approximating this one coherent state transfers the signed
expectation needed below.  No success claim for $\widetilde U$ on
individual basis inputs is required.

Because $\widetilde U$ acts trivially on $R$, and a CNOT conjugates
$Z$ on its target to the product of $Z$ on control and target,
\begin{equation}\label{eq:reference-conjugation}
 W^{\dagger}Z_RW
 =\operatorname{COPY}^{\dagger}Z_R\operatorname{COPY}
 =Z_I Z_R,
 \qquad Z_I=\prod_{i=1}^n Z_{I_i}.
\end{equation}
The one-pair conjugation identity follows directly on a basis vector:
the target bit after copying is $r_i\mathbin{\oplus}x_i$, so its
$Z$ eigenvalue is $(-1)^{r_i+x_i}$.  Acting on $\psi_0$, the last
operator in \eqref{eq:reference-conjugation} gives
\[
 \psi'_0=W^{\dagger}Z_RW\psi_0
 =|-\rangle_I^{\otimes n}
       \otimes|0\rangle^{\otimes a}
       \otimes|0\rangle_R^{\otimes n}.
\]
Let $H_0$ be the count associated with the full product state
$\psi_0$.  Every input factor of $\psi'_0$ is orthogonal to its
specified $|+\rangle$ state, while all other factors agree with their
specified states.  Consequently
\[
 \norm{\psi'_0}=1,\qquad H_0\psi'_0=n\psi'_0.
\]
Moving the unitary and $Z_R$ across the inner product gives
\begin{equation}\label{eq:parity-pairing}
 \ip{\Omega}{Z_RP_B\Omega}
 =\ip{\psi'_0}{W^{\dagger}P_BW\psi_0}.
\end{equation}
This identity does not require commuting $P_B$ through $W$ or
preserving an input register.

Choose fixed exponents
\[
 0<p<h<\min\{1,1/k\}.
\]
By \eqref{eq:parity-resources}, $m^h=O(n^{kh})=o(n)$, so
$m^h<n$ eventually.  The vector $\psi'_0$ then belongs to the range
of $\one_{H_0\ge m^h}$.  Cauchy--Schwarz and
Corollary~\ref{cor:product-tail}, applied to the regular trajectory
of $W$, show that
\[
 \left|\ip{\Omega}{Z_RP_B\Omega}\right|
 \le\norm{\one_{H_0\ge m^h}W^{\dagger}P_BW\psi_0}
 \le\exp(-m^p).
\]
This tends to zero, contradicting \eqref{eq:pruned-correlation}.

For the fixed $d,k,C,\varepsilon$, let $\mathcal S$ be the set of input lengths
admitting an admissible circuit with success at least $1/2+\varepsilon$ on every
input.  The contradiction excludes an unbounded $\mathcal S$;
hence $\mathcal S$ is finite.  There is therefore a threshold
$n_0=n_0(d,k,C,\varepsilon)$ such that at every $n\ge n_0$, every circuit meeting
the specified depth and total-qubit bounds has an input on which
its success probability is strictly less than $1/2+\varepsilon$. In particular,
for every positive integer $N$, every $n\ge\max\{N,n_0\}$ has
this failure property.  This proves
Theorem~\ref{thm:parity} with the stated strict failure threshold.
\end{proof}

\paragraph{Majority and symmetric-function consequences.}
Xu and Li's Theorem~1 and Corollaries~2--3
\citep{XuLi2026Fanout}, combined with Theorem~\ref{thm:parity}, also
exclude strict majority and broader symmetric families in our circuit
model. More precisely, let \(f_n:\{0,1\}^n\to\{0,1\}\) be symmetric, let \(f_{n,i}\) be its
value at Hamming weight \(i\), and set
\(\rho(f_n)=\max_{0\le i<n:\,f_{n,i}\ne f_{n,i+1}}\min\{i+1,n-i\}\), with zero
for constant functions. Its pointwise-error-\(1/3\) approximate degree is
the least total degree of a real polynomial \(p\) satisfying
\(\lvert p(x)-f_n(x)\rvert\le1/3\) for every \(x\in\{0,1\}^n\).
Assume either \(\rho(f_n)\ge n^\delta\) for all sufficiently large \(n\)
for some fixed \(\delta>0\), or this approximate degree is
\(\Omega(n^{1/2+\eta})\) for some fixed \(\eta>0\). Then no
fixed-depth, polynomial-total-qubit family achieves success at least
\(1/2+1/(\log n)^\gamma\) on every input at every sufficiently large
length, for any fixed \(\gamma>0\). Strict majority
\(\operatorname{MAJ}_n(x)=\mathbf1_{\{|x|>n/2\}}\) has
\(\rho=\lceil n/2\rceil\).

The full formulation and output-model adjustment are recorded in the
companion article
\citep[Section~7.3, Corollary~10 and its proof]{OpenAILocalization2026}.
Xu and Li require the measured output to be an initially zero ancilla.
Appending one fresh zero qubit and a CNOT from our designated output
preserves its measurement distribution even with arbitrary garbage, at a
cost of one qubit and one layer.
A family with the stated success would, after this adjustment, yield exact
clean fanout through the cited reduction. The Hadamard-to-parity argument in
Section~\ref{sec:introduction} excludes such fanout using our own
Theorem~\ref{thm:parity}. Under the same circuit conventions, adjoining
unit-cost primitive reversible unbounded threshold gates therefore
strictly enlarges the class computed with worst-case success at least
\(2/3\): the reversible gate
\(\lvert x,b\rangle\mapsto\lvert x,b\mathbin\oplus\operatorname{MAJ}_n(x)\rangle\)
computes majority exactly from a zero target.

\paragraph{State-preparation consequences.}
We justify the bounds stated in Section~\ref{sec:introduction}.
Keep every register of a preparation when using its unitary and inverse.
For felinity $\eta>0$ on $h$ retained qubits of a depth-$d$ preparation
on $M$ total qubits, the constructions in the proofs of
\citet[Lemmas~4.9 and~4.1]{GGJ2026Fourier} give target probabilities
$a\ge1/2$ on $0^h$ and $b\ge\eta^2/2$ on $1^h$, in depth $O(d+1)$
on $O(M+h)$ qubits. The square occurs because the top Fourier
coefficient has magnitude at least $\eta$, so its weight is at least $\eta^2$.

The exact step is the one-round case of
\citet[Theorem~7]{GrierMorrisWu2026}. For unit vectors $g,e$, if a
unitary $V$ prepares
$V|0\rangle=(|g,0\rangle+\sqrt3|e,1\rangle)/2$, where $|0\rangle$
denotes its entire zero input and the last qubit is a flag, then
\[
 V(I-2|0\rangle\langle0|)V^\dagger Z_{\rm flag}V|0\rangle
 =|g,0\rangle.
\]
Indeed, $\langle0|V^\dagger Z_{\rm flag}V|0\rangle=-1/2$.
A flag initialized to one can be given coefficient
$c=1/(2\sqrt{a+b})\le1$ on zero separately under the two extreme
target tests: use a one-qubit reflection sending $|1\rangle$ to
$c|0\rangle+\sqrt{1-c^2}|1\rangle$. Each conditional operation is a
product reflection, so the two tests use two sequential Toffoli gates
up to local conjugations. The flag-zero component is one half of the
normalized union of the extreme branches. The displayed identity
prepares that union exactly, with one-branch probability
$q=b/(a+b)\in[\eta^2/2,1/2]$, since $a+b\le1$ and $b\le1-a\le a$.

Prepare $r=\lceil(\log 2)/(-\log(1-q))\rceil=O(\eta^{-2})$
independent copies on fresh zeros. Disjoint columnwise reversible ORs
into $h$ new targets give probabilities
$u=(1-q)^r\in[1/4,1/2]$ on $0^h$ and $v=1-u\in[1/2,3/4]$ on $1^h$.
Here $(1-q)/2\le(1-q)^r\le1/2$ and $r\le(\log 2)/q+1$.
Giving these two branches flag-zero coefficients
$1/(2\sqrt{2u})$ and $1/(2\sqrt{2v})$, respectively, and applying
the same exact round yields
\[
 2^{-1/2}\bigl(|0^h,\alpha_0\rangle+|1^h,\alpha_1\rangle\bigr),
\]
with unit branch states $\alpha_0,\alpha_1$ containing arbitrary
garbage. The all-zero reflection is likewise a local conjugate of a Toffoli,
and the ORs use Toffoli gates with local $X$ gates. The source-dependent
flag coefficients are permitted by nonuniformity. Rosenthal's exact
case \citep[Theorem~3.1]{Rosenthal2021} converts this balanced state,
using its complete preparation and inverse, to the exact parity unitary
on arbitrary states of the $h+1$ data qubits, with every auxiliary
qubit restored to zero, and hence to clean $F_h$ by Hadamard conjugation. All copies
run in parallel, and only a constant number of preparation and
reflection stages are sequential. The resulting depth is $O(d+1)$
and the total-qubit bound is $O((M+h)\eta^{-2}+h)$.

For the Dicke bound, use the block calculation in the proof of the
companion's state-preparation corollary
\citep[Section~7.1, Corollary~8 and its proof]{OpenAILocalization2026}.
Uniformly for $n\ge2k$ and all sufficiently large $k$, it shows that
an input at trace distance at most $1/(80k)$ from $|D_k^n\rangle$
produces, with $\ell=\lfloor k/\log k\rfloor$ fresh targets and three
additional allowed layers, a retained state of felinity greater than
$1/(40k)$. Applying the preceding promotion at arity $\ell$ gives
exact parity in fixed depth on $O((M+n)k^2)$ total qubits. For fixed
$0<\delta<1$ and $k\ge n^\delta$, eventually $\ell\ge\sqrt k$, so
$n\le k^{1/\delta}\le\ell^{2/\delta}$ and $k^2\le\ell^4$ make this one fixed polynomial
bound in $\ell$. A felinity violation $\eta\ge n^{-A}$ similarly
preserves one fixed polynomial bound at arity $n$. Both parity arities
grow along any unbounded sequence of violations, contradicting the
eventual pointwise Theorem~\ref{thm:parity}.
